\subsection{群}

回顾以下定义（§2，第3号，定义6）。

定义1. 一个具有结合律的合成法则、含有单位元且其中每个元素都可逆的集合，称为群。

换言之，群是其中每个元素都可逆的幺半群（§2，第1号，定义1）。集合上确定其群结构的合成律称为群律。若G和H是两个群，则从G到H的岩浆同态也称为群同态。这样的同态f将单位元映至单位元；事实上，设e（分别为e'）为G（分别为H）的单位元；将G和H的群律写成乘法形式，有e·e = e，从而 \(f(e).f(e) = f(e)\) 且乘以 \(f(e)^{-1}, f(e) = e'\) 因此 f 是含幺的。于是由第 2 节第 3 条可知 \(f(x^{-1}) = f(x)^{-1}\) 为所有 \(x \in G\) .

例。在任何幺半群 E 中，可逆元素集在 E 上诱导的结构下构成一个群。特别地，集合 F 到自身的双射映射集（或 F 的置换集）在运算 \((f, g) \mapsto f \circ g\) 下构成一个群，称为集合 F 的对称群，记作 \(S_{F}\) .

在本段中，除非另有说明，群的合成法则将始终以乘法形式书写，且e将表示该群法则的单位元。

群G称为有限群，如果G的底集是有限的；否则称为无限群；群的基数称为群的阶。

若G上的一个合成法则确定了G上的群结构，其对偶法则亦然。将群G映到自身的映射，将每个元素对应到 \(x \in G\) x 的逆是 G 到其相反群的一个同构。 \(\S 2\) ，第3号，命题4）。

按照我们的一般约定（《集合论》第二卷，第3节，第1号），我们将用 \(\mathbf{A}^{-1}\) 表示映射下G的子集A的像。 \(x \mapsto x^{-1}\) 但重要的是要注意，尽管记号类似， \(\mathbf{A}^{-1}\) 绝不是A在G的子集之间的合成法则下的逆元素（回忆一下，XY是满足 \((\mathbf{X}, \mathbf{Y}) \mapsto \mathbf{XY}\) 的xy的集合）：该法则下的单位元是 \(x \in \mathbf{X}, y \in \mathbf{Y}\) ，而在此法则下 \(\{e\}\) 中仅有的可逆元素是由单个元素组成的集合A（这样的A，当然，其逆元素确实是 \(\mathfrak{P}(\mathbf{G})\) ）。单位元 \(\mathbf{A}^{-1}\) 。A称为 \((\mathbf{AB})^{-1} = \mathbf{B}^{-1}\mathbf{A}^{-1}\) 对……成立 \(\mathbf{A} \subset \mathbf{G}, \mathbf{B} \subset \mathbf{G}\) 的对称子集， \(\mathbf{G}\) 中正规，若 \(\mathbf{A} = \mathbf{A}^{-1}\) 的分拆集合的双射。对于所有 \(\mathbf{A} \subset \mathbf{G}, \mathbf{A} \cup \mathbf{A}^{-1}, \mathbf{A} \cap \mathbf{A}^{-1}\) 和 \(\mathbf{AA}^{-1}\) 是对称的。

